\scanpage{280}
% archive foliation: 279

\unsure{Cette} Commission comprend tous les maîtres de conférences et professeurs du Laboratoire de Mathématique qui formeront au moins deux tiers de ladite Commission. Celle ci pourra inclure également, le cas échéant, des chercheurs détachés du CNRS au laboratoire de Math. et des Maître de conférence et professeurs de l'Ecole de Mathématique ayant compétence pour juger le candidat. La nomination est faite sur avis favorable de deux tiers au moins des membres de la Commission.

IV Détachements et changements de Sections.

1°/ Les membres du personnel enseignant de l'Ecole de Mathématique ont droit, tous les 4 ans, à un détachement auprès du Laboratoire, pour une durée d'un semestre. Ils n'ont aucune obligation d'enseignement pendant cette période.

Sur leur demande, et avec l'accord de l'Assemblée du Département, ils peuvent être détachés auprès du laboratoire pour une période d'un an \unsure{renouvelable}. Cette possibilité sera offerte, en \unsure{particulier}, aux enseignants, Docteurs d'Enseignement ou \unsure{non}, désirant obtenir la qualification nécessaire pour faire partie du laboratoire à titre \unsure{permanent}.

2°/ Les membres du personnel enseignant du Laboratoire \unsure{peuvent}, sur leur demande, et avec l'accord de l'\unsure{Assemblée} du Département, être détachés auprès de l'Ecole pour une période d'un an renouvelable.
